\subsubsection{Unabhängige
Rekonstruktionsprüfung}\label{unabhuxe4ngige-rekonstruktionspruxfcfung}

Die Familie wird ohne Zustandsdaten festgelegt. Anschließend werden die
beiden obigen reinen Zustände sowie sechs deterministisch zufällige
Zustände der Ränge 1, 2, 5, 17, 60, 60 in derselben ursprünglichen
W-Codebasis verwendet. Aus ihren 3599 kalibrierten Schatten
rekonstruiert der Prüfer die vollständigen Dichtematrizen. Zusätzlich
sagt er je zwölf vorher unbenutzte, G-gedrehte Paarerwartungswerte
dieser selben Zustände voraus.

Diese Tests sind \textbf{numerisch}, nicht als exakte Zustandsbeweise
ausgegeben. Im dokumentierten Lauf beträgt der größte
Matrixeintragsfehler unter 7.6×10⁻¹⁶; der größte Fehler der 96
unbenutzten Paarvorhersagen liegt unter 5.8×10⁻¹⁷. Die Antwortmatrizen
der Faktoren haben Konditionszahlen ungefähr 40.47 und 20.72. Das ist
ein rauschfreier Rekonstruktionstest; statistische Messkosten,
Zustandspräparation und Robustheit gegen reales Rauschen sind hier nicht
nachgewiesen.

\subsection{5. Was native Zeitentwicklung
hinzufügt}\label{was-native-zeitentwicklung-hinzufuxfcgt}

Der gemeinsame helle N=2-Prozess besitzt H=h⊗I\_60. Für jeden
Zweigoperator B und inneren Operator O gilt exakt

\[\adjustbox{max width=\linewidth}{$\displaystyle 
 [h\otimes I,B\otimes O]=[h,B]\otimes O.
$}\]

Beliebige Zeitableitungen und Zeiten verändern daher nur den
Zweigfaktor. Sie erweitern den inneren Schattenrang 24 beziehungsweise
736 nicht. Die zusätzlichen Paarsonden können den inneren Rang dagegen
bis 3600 erhöhen. Zweigkohärenz und ihre Rekonstruktion erfordern eine
gesonderte Prozessbetrachtung; eine innere Zustandsrekonstruktion allein
bestimmt noch keine vollständige Mehrzeitgeschichte oder ein
Messinstrument.

\subsection{Reproduktion und Grenze}\label{reproduktion-und-grenze}

Mit Python 3.10 oder neuer und NumPy, aus diesem Ordner:

\begin{verbatim}
python3 -B replay.py
\end{verbatim}

Der Lauf überprüft die fünf eingefrorenen Quellenpins, führt den eigenen
Prüfer normal und mit \codeword{-OO} aus und verlangt byte-identische
Ausgaben. Warnungen gelten als Fehler; keine Prüfbedingung benutzt
abschaltbare Python-Assertions. Die Prüfzahlen stehen im
Replay-Manifest. Die vier Kontextdateien werden gehasht, nicht als
fremde Gesamtsuite erneut ausgeführt. Die mathematische Prüfung benutzt
nur die lokale W-Archivdatei.

Damit ist eine gemeinsame, am Ursprungstensor kalibrierte \emph{innere}
Schattenkarte konstruiert. Ein räumlicher Rand, ursprüngliche
ausführbare Instrumente, eine universelle Prozessidentifikation,
3+1D-Raumzeit oder ein Abschluss von T1-T8 folgt daraus nicht.

\begin{center}\rule{0.5\linewidth}{0.5pt}\end{center}

\section{Teilbericht C: kausaler
Eingriff}\label{teilbericht-c-kausaler-eingriff}

\section{Exakter kausaler Eingriff am unveränderten nativen
Hamiltonoperator}\label{exakter-kausaler-eingriff-am-unveruxe4nderten-nativen-hamiltonoperator}

\subsection{Ergebnis}\label{ergebnis-2}

Im ursprünglichen TFPT-Modell gibt es einen vollständig gelösten
kausalen Test innerhalb des Ladungssektors N=2. Ein Phasenimpuls an
Fermionmode 4 ändert die spätere unbedingte Besetzung von Mode 5. Die
beiden Observablen kommutieren vor der Entwicklung. Der Impuls erhält
die Ladung und wirkt nur auf Mode 4; es gibt keine Konditionierung auf
Messergebnisse.

Der stärkste hier bewiesene Vergleich endet in \textbf{beiden} Armen mit
genau zwei Fermionen und keinem Boson. Gemessen wird eine andere
Verteilung der Fermionen auf disjunkte Paare. Die natürliche Entwicklung
zwischen Eingriff und Messung verwendet ausschließlich das ursprüngliche
H und W.

\textbf{Bedingung:} Die Präparation des angegebenen Produktzustands, ein
gezielter Phasenimpuls und die einzelne Besetzungsmessung werden als
Instrumente gewährt. Ihre physische Verfügbarkeit aus dem ursprünglichen
Operationssatz ist hier nicht hergeleitet. Der Impuls ist eine
ausdrücklich hinzugewährte Kontrolloperation; er wird nicht als von der
ungestörten H-Entwicklung bereitgestellt ausgegeben. Es wurde kein
zusätzlicher Hopping- oder Linkterm in die freie Dynamik eingeführt.

Das Ergebnis betrifft Moden innerhalb derselben Bank. Es enthält weder
eine räumliche Metrik noch den bereits untersuchten N=64-Grundzustand
oder dessen geladenen Entnahmepol.

\subsection{1. Quelle und vollständiger invarianter
Sektor}\label{quelle-und-vollstuxe4ndiger-invarianter-sektor}

Die eingefrorene Quelle ist \codeword{sources/spinor\_tensors.npz} mit
SHA-256
\codeword{3f00a0892157a691e4ef26920c3702772c7bf76a250fad98ff04a079bb64b763}.

Es gilt unverändert

\[\adjustbox{max width=\linewidth}{$\displaystyle 
H=\Delta N_b+g\sum_A(b_A^\dagger P_A+P_A^\dagger b_A),\qquad
P_A=\sum_{i<j}W_{A,ij}f_jf_i,\qquad N=N_f+2N_b.
$}\]

W besitzt 60 Zeilen, 2016 Spalten und 480 Einträge mit Wert ±1. Jede
aktive Spalte gehört genau einer Zeile. Jede Zeile enthält acht
Fermionenpaare mit insgesamt sechzehn verschiedenen Moden. Deshalb
zerfällt der \textbf{gesamte} N=2-Sektor in 60 invariante
neundimensionale Blöcke und 1536 vollständig dunkle Paarzustände. Die
sieben dunklen Richtungen innerhalb jedes Blocks kommen zusätzlich
hinzu: insgesamt bleibt die bekannte dunkle Dimension 1536+60·7=1956.

Eine solche Gruppe besteht aus einem Boson und acht Paaren. Mit
\(|p_\ell\rangle=f_{i_\ell}^\dagger f_{j_\ell}^\dagger|0\rangle\),
\(s_\ell=W_{A,i_\ell j_\ell}\) und
\(|S\rangle=\sum_\ell s_\ell|p_\ell\rangle/\sqrt8\) gilt

\[\adjustbox{max width=\linewidth}{$\displaystyle 
H_A=\begin{pmatrix}0_{8\times8}&g s\\g s^T&\Delta\end{pmatrix},
\qquad H_A|_{\operatorname{span}(S,b_A)}=
\begin{pmatrix}0&\sqrt8g\\\sqrt8g&\Delta\end{pmatrix}.
$}\]

Die CAR-Vorzeichen wurden für alle 480 aktiven Paare separat geprüft.
Dies ist eine exakte Einschränkung auf einen invarianten Sektor, keine
Abschneidung höherer Bosonstufen.

Die benutzte Zeile A=0 lautet:

{\def\LTcaptype{none} % do not increment counter
\begin{longtable}[]{@{}rlr@{}}
\toprule\noalign{}
Blatt & Fermionenpaar & Vorzeichen \\
\midrule\noalign{}
\endhead
\bottomrule\noalign{}
\endlastfoot
0 & (4,57) & −1 \\
1 & (5,56) & +1 \\
2 & (8,53) & +1 \\
3 & (9,52) & −1 \\
4 & (16,45) & −1 \\
5 & (17,44) & +1 \\
6 & (28,33) & +1 \\
7 & (29,32) & −1 \\
\end{longtable}
}

\subsection{2. Lokale Algebra und ehrlicher
Interventionstest}\label{lokale-algebra-und-ehrlicher-interventionstest}

Die Senderoperation ist

\[\adjustbox{max width=\linewidth}{$\displaystyle 
Z_4=(-1)^{n_4}=I-2n_4,\qquad
\mathcal E_4(\rho)=Z_4\rho Z_4^\dagger.
$}\]

Sie hat genau einen unitären Krausoperator und ist daher vollständig
positiv und spurerhaltend. Sie ist fermionparitätsgerade und erhält N.
Auf der ganzen Fock-Algebra kommutiert sie mit allen disjunkten
Modenoperatoren; insbesondere \([Z_4,n_5]=0\). Ihre Einschränkung auf
den ausgewählten Block wechselt allein das Vorzeichen des Blattes
(4,57). Diese Einschränkung wurde gegen den wirklichen Operator
\((-1)^{n_4}\) auf allen 2076 N=2-Basiszuständen geprüft.

Der Empfänger liest \(B=n_5\) aus. Im gewählten invarianten Block ist
dies der Projektor auf Blatt (5,56). Die Operation enthält keinen
Verweis auf Mode 5. Zu ihrem Zeitpunkt bleibt dessen gesamte
Besetzungsstatistik für jeden Zustand identisch, weil
\(Z_4^\dagger n_5 Z_4=n_5\).

Mit \(\rho_T=U_T|p_0\rangle\langle p_0|U_T^\dagger\) ist die gemessene
Differenz genau

\[\adjustbox{max width=\linewidth}{$\displaystyle 
\delta_{4\to5}(T)=
\operatorname{tr}\!\left[n_5U_T\mathcal E_4(\rho_T)U_T^\dagger\right]
-\operatorname{tr}\!\left[n_5U_T\rho_TU_T^\dagger\right].
$}\]

Das ist ein Vergleich von unbedingten Wahrscheinlichkeiten nach einem
kontrollierten lokalen Eingriff. Eine Zweipunktkorrelation wird dafür
nicht als Ersatz benutzt. Die vor dem Impuls erzeugte Kohärenz wird
ausdrücklich durch die native Vorentwicklung präpariert. Ein
Phasenimpuls direkt auf dem anfänglichen Produktzustand hätte nur ein
globales Vorzeichen und überhaupt keinen Effekt; diese Negativkontrolle
wird ebenfalls geprüft.

\subsection{3. Exakter Test mit gleicher
Endzusammensetzung}\label{exakter-test-mit-gleicher-endzusammensetzung}

Setze für \(\Delta>0\) und reelles \(g\ne0\)

\[\adjustbox{max width=\linewidth}{$\displaystyle 
\Omega=\sqrt{\Delta^2/4+8g^2},\quad
T=\frac\pi\Omega,\quad d=\frac\Delta{2\Omega},\quad
z=-e^{-i\pi d},\quad
x=|z-1|^2=4\cos^2\frac{\pi d}{2}.
$}\]

Nach T wirkt der helle Zweizustandsblock als z mal Identität; die
dunklen Paarzustände bleiben unverändert. Es gibt dann für \textbf{jede}
anfängliche Paarkombination keine Bosonamplitude. Auf der Paarseite ist

\[\adjustbox{max width=\linewidth}{$\displaystyle 
U_T^{(f)}=I_8+(z-1)|S\rangle\langle S|.
$}\]

Beide Arme beginnen in demselben vollständig angegebenen Produktzustand
\(|p_0\rangle=f_4^\dagger f_{57}^\dagger|0\rangle\), alle anderen Moden
leer. Sie unterscheiden sich nur durch den mittleren Phasenimpuls:

\[\adjustbox{max width=\linewidth}{$\displaystyle 
\begin{array}{lll}
\text{unberührt:}&U_T\;I\;U_T|p_0\rangle,&\text{danach }n_5,\\
\text{Eingriff:}&U_T\;Z_4\;U_T|p_0\rangle,&\text{danach }n_5.
\end{array}
$}\]

Für jedes andere Blatt q sind die exakten Endamplituden

\[\adjustbox{max width=\linewidth}{$\displaystyle 
\langle p_q|U_T^2|p_0\rangle
=s_0s_q\frac{z^2-1}{8},\qquad
\langle p_q|U_TZ_4U_T|p_0\rangle
=s_0s_q\frac{3(z-1)^2}{32}.
$}\]

Damit lautet insbesondere die unbedingte Empfängerwahrscheinlichkeit

\[\adjustbox{max width=\linewidth}{$\displaystyle 
p_{\rm frei}(n_5=1)=\frac{x(4-x)}{64},\qquad
p_{\rm Impuls}(n_5=1)=\frac{9x^2}{1024},
$}\]

und die kausale Differenz ist

\[\adjustbox{max width=\linewidth}{$\displaystyle 
\boxed{\delta_{4\to5}(T)=\frac{x(25x-64)}{1024}}.
$}\]

Für \(0<g^2\le3\Delta^2/32\) gilt \(1/2\le d<1\), also \(0<x\le2\) und
folglich \textbf{strikt \(\delta_{4\to5}<0\)}. Das ist eine exakte
endliche Zeit, kein unkontrollierter Schluss aus einer numerisch kleinen
Antwort oder einer asymptotischen Reihe. Außerhalb dieses Intervalls
gilt weiterhin die exakte Formel, kann aber auch das Vorzeichen wechseln
oder an x=64/25 verschwinden.

Am bisherigen Prüfpunkt \(\Delta=1,g=1/20\) ist \(d=5/(3\sqrt3)\) und
\(T\approx6.0459978807807\):

{\def\LTcaptype{none} % do not increment counter
\begin{longtable}[]{@{}lr@{}}
\toprule\noalign{}
Größe & Wert \\
\midrule\noalign{}
\endhead
\bottomrule\noalign{}
\endlastfoot
Ohne mittleren Impuls & 0.00087491599417248 \\
Mit mittlerem Impuls & 0.0000017344871305812 \\
Differenz & −0.00087318150704190 \\
Endbesetzung in beiden Armen & Nf=2, Nb=0 exakt \\
\end{longtable}
}

Die Wellenfunktion kann zwischen den Ablesezeiten einen Boson enthalten.
Das natürliche H vermittelt dadurch die Umverteilung von einem
Fermionenpaar in andere. Die Endmessung benötigt keine Selektion auf
„kein Boson``: Dieser Zustand gilt in beiden Armen mit
Wahrscheinlichkeit eins.

\subsection{4. Zwei ergänzende
Kontrollen}\label{zwei-erguxe4nzende-kontrollen}

\subsubsection{Boson als einfachere
Referenz}\label{boson-als-einfachere-referenz}

Ausgehend vom Produktzustand \(b_0^\dagger|0\rangle\) wähle die beiden
Verzögerungen \(\tau=\pi/(2\Omega)\). Der unberührte Arm kehrt exakt zum
Boson zurück, also \(p_{\rm frei}(n_5=1)=0\). Ein mittleres Z4 erzeugt

\[\adjustbox{max width=\linewidth}{$\displaystyle 
p_{\rm Impuls}(n_5=1)
=\frac{(1-d^2)\left[1+d^2-2d\sin(\pi d/2)\right]}{128}
\ge\frac{(1-d^2)(1-d)^2}{128}>0.
$}\]

Dies gilt für jedes \(g\ne0\) und \(\Delta>0\). Auch dieser Test ist
ladungserhaltend und unbedingt. Er hat allerdings unterschiedliche
Fermion/Boson-Zusammensetzungen am Ende; für die stärkere Aussage über
reine Paarumverteilung wird deshalb der Test aus Abschnitt 3 benutzt.

\subsubsection{Lokales Auffüllen und kurze
Zeiten}\label{lokales-auffuxfcllen-und-kurze-zeiten}

Auf der Referenz \(f_{57}^\dagger|0\rangle\) füllt das lokale
CPTP-Instrument mit Krausoperatoren \(n_4,f_4^\dagger\) Mode 4. Ohne
Eingriff ist der Einfermionzustand stationär. Mit Eingriff entsteht
\(|p_0\rangle\), woraus

\[\adjustbox{max width=\linewidth}{$\displaystyle 
\delta n_5(t)=\frac{|a(t)-1|^2}{64},\qquad
a(t)=e^{-i\Delta t/2}\left[\cos\Omega t+i d\sin\Omega t\right].
$}\]

Die führenden Beiträge sind \(\delta n_5(t)=g^4t^4/4+O(t^6)\) und
\(\delta N_{b_0}(t)=g^2t^2+O(t^4)\). Der Pfad in ein anderes
Fermionenpaar hat somit zwei native H-Anwendungen in der Amplitude; die
erste Umwandlung zum Boson hat eine.

Dieses Auffüllen tauscht eine Ladungseinheit mit einem angenommenen
Reservoir aus. Es wird nicht als zahlenerhaltende Operation ausgegeben
und ist für den Hauptnachweis nicht erforderlich.

\subsection{5. Prüfung und unveränderte offene
Grenze}\label{pruxfcfung-und-unveruxe4nderte-offene-grenze}

\codeword{verify\_causal.py} benutzt explizite Fehlerbedingungen, keine
abschaltbaren Assertions. Normaler Python-Lauf und
\codeword{python -OO}, jeweils mit Warnungen als Fehler, erzeugen
byte-identische JSON-Ergebnisse. Der optimierte Lauf wurde aus einem
anderen Arbeitsverzeichnis gestartet; die Quelle wird relativ zum
Prüfprogramm aufgelöst. Abhängigkeiten: NumPy, SymPy und SciPy.

Die symbolischen Identitäten und die vollständige W-Struktur werden
exakt geprüft. Zusätzlich läuft ein unabhängiger numerischer Vergleich
auf dem vollen 2076-dimensionalen N=2-Hamiltonoperator, direkt aus W
aufgebaut. Er bestätigt beide vollständigen Endzustände, die
Wahrscheinlichkeiten, Normerhaltung, verschwindende Bosonamplituden und
die Negativkontrolle. Der abschließende Lauf enthält \textbf{662 exakte
und 10 numerische Prüfbedingungen}. Prüfzahlen und Quellhash stehen in
\codeword{result.json}; \codeword{result\_OO.json} ist dessen
identischer optimierter Replay. Beide Ergebnisdateien haben SHA-256
\codeword{b37ae758117e4c35d1063e76756cd6a39f133c964d331a9d400e67a0b157a4a6}.

\textbf{Geschlossen:} Im ausdrücklich gewährten Instrumentenvertrag
existiert am selben ursprünglichen H ein kausaler
Moden-zu-Moden-Eingriff. Ein weiteres Hopping-Glied ist für diesen
N=2-Zeugen mathematisch nicht notwendig. Reine Paarumverteilung wird am
Endzeitpunkt getrennt von Speziesumwandlung bewiesen.

\textbf{Weiter offen:} Auswahl und Präparation dieses
N=2-Referenzzustands, ausführbare einzelne Modenphasen und
Besetzungsinstrumente aus dem nativen Operationssatz; Anschluss an den
N=64-Grundzustand und dessen geladenen Pol; Interpretation der
Modenalgebren als räumlich unabhängige Labore; Distanz, Raumdimension,
endliche Ausbreitungsgeschwindigkeit und gemeinsame Raumzeit. Die bloße
Tatsache, dass ein Kontrolloperator zur mathematischen Fock-Algebra
gehört, schließt diese Operationsfrage nicht.

\begin{center}\rule{0.5\linewidth}{0.5pt}\end{center}

\section{Teilbericht D: derselbe Dreizustandsraum für Schatten und
Eingriff}\label{teilbericht-d-derselbe-dreizustandsraum-fuxfcr-schatten-und-eingriff}

\section{Derselbe Dreizustandsraum trägt kausalen Eingriff und
Zustandstomografie}\label{derselbe-dreizustandsraum-truxe4gt-kausalen-eingriff-und-zustandstomografie}

